\chapter{Future Work}\doublespacing

\section{Overview}
While the current implementation of the \textbf{LightAngo} compiler establishes a fully functional, robust compilation pipeline with mandatory documentation enforcement and native x86-64 code generation, several promising avenues for future research and technical enhancement have been identified.

Future developments will focus on expanding the language's expressive power, deepening semantic documentation verification, integrating industry-standard intermediate representation backends, and enhancing developer tooling.

\section{Future Enhancements}

\subsection{1. Advanced Data Structures and Type System Expansion}
The current version of LightAngo primarily supports scalar 64-bit integer values, expressions, and string literals. Future iterations will incorporate:
\begin{itemize}
    \item \textbf{Compound Data Types:} Multi-dimensional arrays, user-defined structures (`struct`), and algebraic data types.
    \item \textbf{Explicit Typing and Type Inference:} Support for static type annotations (`let x: int = 10;`, `let f: float = 3.14;`) combined with Hindley-Milner type inference.
    \item \textbf{Heap Memory Management:} Dynamic memory allocation primitives (`alloc`, `free`) coupled with scoped RAII-based automatic cleanup or reference counting.
\end{itemize}

\subsection{2. Natural Language Processing (NLP) Documentation Quality Assessment}
Currently, the comment validation mechanism enforces lexical existence and minimum character length metrics. An advanced future enhancement will integrate lightweight Natural Language Processing (NLP) heuristics and embedding models to assess:
\begin{itemize}
    \item \textbf{Semantic Relevance:} Measuring cosine similarity between function parameter names, variable identifiers, and words utilized in the doc-comment.
    \item \textbf{Completeness of Contract:} Verifying that documented pre-conditions, post-conditions, and loop invariants mathematically align with AST branch conditions.
\end{itemize}

\subsection{3. Register Allocation via Graph Coloring}
The current code generator relies primarily on a stack-based memory layout for local variables and intermediate expressions. Future work will implement an advanced machine-dependent optimizer featuring:
\begin{itemize}
    \item \textbf{Liveness Analysis:} Constructing interference graphs representing overlapping variable lifetimes.
    \item \textbf{Kempe/Chaitin Graph Coloring:} Mapping virtual variable nodes to physical hardware general-purpose registers (`%rax`, `%rbx`, `%rcx`, `%rdx`, `%rsi`, `%rdi`, `%r8`--`%r15`), minimizing memory stack spill operations by over 80\%.
\end{itemize}

\subsection{4. Multi-Architecture Backend via LLVM IR Integration}
To enable cross-platform binary synthesis across heterogeneous computing environments, a future module will emit typed LLVM Intermediate Representation (LLVM IR):
\begin{itemize}
    \item Enabling compilation targeting ARM64 (Apple Silicon, Raspberry Pi), RISC-V, WebAssembly (WASM), and Windows PE x86-64.
    \item Leveraging LLVM's industrial-grade optimization pipeline (O1, O2, O3, LTO, vectorization).
\end{itemize}

\subsection{5. Language Server Protocol (LSP) and IDE Extensions}
To enhance developer ergonomics, future work will encapsulate the compiler frontend into an asynchronous Language Server conforming to the Language Server Protocol (LSP):
\begin{itemize}
    \item Providing real-time documentation warning squiggles, auto-completion of doc-comment templates, and jump-to-definition directly within VS Code, CLion, and Neovim.
\end{itemize}

\subsection{6. Visual AST and Execution Tracing Tools}
Developing an interactive graphical web interface to render AST trees, symbol table scope transitions, and step-by-step register states during binary execution to maximize LightAngo's educational utility in university compiler construction curricula.
\newpage
